% Page 048

\seite{48}

\abschnitt{Weihnachtsfeierlichkeiten}
\pstart
Der Geburtstag Sr. Majestät des Kaisers wurde von den Kindern der Pfarrei\ob
gemeinsam im Pfarrsaale zu Seffern unter Leitung des Herrn Ortsschul\ohb
inspektors gefeiert. Vaterlandslieder, die von den größeren Kindern\ob
gesungen, wurden Gedichte vorgetragen. Die Ansprache an die Kinder\ob
hielt Herr Pfarrer aus Seffern.

\pend
\abschnitt{Sefferweich 31. März 1911}
\pstart


\pend
\begin{verse}
Zimmerer Th.
\end{verse}
\pstart

\pend
\abschnitt{Osterprüfung}
\pstart

Die Osterprüfung wurde am 31.\ März durch den Ortsschul\ohb
inspektor Herrn Pfarrer Zimmerer aus Seffern abgehalten.

\pend
\jahresueberschrift{Das Schuljahr 1911-12}
\pstart

Entlassen wurden im Herbste dieses Jahres 6 Kinder, und zwar\ob
5 Knaben und 1 Mädchen. 4 Knaben widmen sich der\ob
Landwirthschaft. 1 Knabe und das Mädchen widmen sich dem\unsicher\ob
Schuhmacher\unsicher{} widmen. Der Knabe kam in Aufnahme in die\ob
Post- Prüm. Als Aufnahme\unsicher{} das Mädchen in Pries.\unsicher{} Oberstes\unsicher\ob
von Aufgeboten\unsicher{} wurde an 3 Mädchen. Die Schülerzahl\ob
beträgt 23, davon sind 13 Knaben und 10 Mädchen. Alle Kinder\ob
derselben sind katholisch. Ein Knabe ist nicht willfährig.

\pend
\abschnitt{Seffern 15. 5. 11}
\pstart


\pend
\begin{verse}
{[Unterschrift]}
\end{verse}
\pstart

\pend
\abschnitt{Revision}
\pstart

Die hiesige Schule wurde am 15.\ Mai durch Herrn Kreisschulrat Lenz aus\ob
Bitburg revidiert.

\pend
\abschnitt{Spaziergang}
\pstart

Am 29.\ August unternahmen die vier Schulen der Pfarrei\ob
gemeinschaftlich einen Spaziergang nach dem Kyllthal\unsicher.
\pend
